\scanpage{568}
\mspage{477}
\uline{4.1\deleted{\unsure{5}}6.7.} \uline{Retour sur les sous-catégories fermées d'une
catégorie accessible} (4.16.1).

\uline{Proposition} 4.16.7.1. \unsure{Soient|Soit} $M$ une catégorie
accessible, \deleted{$\pi_0$ un cardinal \inserted{\unsure{infini}} tel que $M$ satisfasse
à $L_{\pi_0}$,} $(\underline S_i)_{i\in I}$ une famille de sous-catégories
\inserted{fermées} accessibles dans $M$\deleted{, et $L_{\pi_0}$-stables
dans $M$}. On ne suppose pas que \unsure{card} $I$ soit
petit. Soit $\underline S$ la sous-catégorie \inserted{str.} pleine
de $M$ telle que $\deleted{\unsure{\underline S}}\operatorname{Ob}\underline S=\bigcup_{i\in I}\operatorname{Ob}\underline S_i$.
Alors pour que $\underline S$ soit accessible dans
$M$, il f. et s. qu'il existe une
\uline{petite} partie $J\subset I$, telle que
$\operatorname{Ob}\underline S=\bigcup_{i\in J}\operatorname{Ob}\underline S_i$.
\inserted{(En particulier, si $I$ est ordonné et $i\mapsto\underline S_i$ \deleted{\unsure{str.}} croissante, \inserted{et \unsure{gd|grd} devant \unsure{tt|ttes} petites parties,} il f. et s. qu'il existe un $i\in I$ tel que \deleted{\unsure{$\underline S_i$}} $\underline S=\underline S_i$, i.e tel que $j\geq i$ implique $\underline S_j=\underline S_i$.)}

C'est évidemment suffisant, en
vertu de 4.13.\ldots\ Prouvons \unsure{inversement|l'assertion} que
si $\underline S$ est accessible \unsure{dans}\inserted{$M$}, alors la condition
précédente est satisfaite. Soit $\pi$ un cardinal
$\geq\pi_0$ tel que $\underline S$ soit $\pi$-accessible dans
$M$, donc $\underline S\xleftarrow{\approx}\operatorname{Ind}_\pi\underline S(\pi)$, où $\underline S(\pi)=\underline S\cap M_\pi$.
% archive sheet number: 567
